\subsection{Fixed stacks for torus actions}\label{rl-000C}
\begin{lemma}\label{rl-000D}
  \label{lem:torus-actions-are-nice}
  \textit{(Torus actions are nice.)}~
  Let \(S\) be a locally noetherian scheme and \(T\) a split torus over \(S\) acting on an algebraic stack \(\mathcal X\) over \(S\).
  \begin{enumerate}[(i)]
    \item If \(\mathcal X\) is Deligne-Mumford, quasi-separated and locally of finite type over \(S\), then \(\eta:\mathcal X^T\to \mathcal X\) is a closed immersion. In particular \(\mathcal X^T\) is algebraic and \(\eta\) is proper and quasi-compact.
    \item If \(\mathcal X\) has affine, finitely presented diagonal, then \(\mathcal X^T\) is algebraic and \(\eta\) is representable by schemes, separated and locally of finite presentation. In particular \(\eta\) is DM-type.
    \item If \(S = \Spec k\) and \(\mathcal X\) is locally of finite type over \(k\) with quasi-compact separated diagonal and affine stabilizers, then \(\mathcal X^T\) is algebraic \cite[Thm. A.20, Rmk A.19]{aranha-khan-latyntsev-etal_VirtualLocalizationRevisited2025} and \(\eta\) is representable by algebraic spaces.
  \end{enumerate}
\end{lemma}
\begin{proof}
$ $
  \begin{enumerate}[(i)]
    \item \cite[Prop A.23]{aranha-khan-latyntsev-etal_VirtualLocalizationRevisited2025}
    \item \cite[Thm 1.2.1 (1)]{romagny_AlgebraicitySmoothnessFixed2022}
          \item The derived stack \(\mathcal X^{hT}\) is 1-Artin by \cite[Thm. A.20]{aranha-khan-latyntsev-etal_VirtualLocalizationRevisited2025}, so its classical truncation \(\mathcal X^T\) is an algebraic stack by Remark \ref{rmk:algebraicity-is-representability} below.
  \end{enumerate}
\end{proof}

\begin{remark}\label{rl-000E}\label{rmk:algebraicity-is-representability}
  The representability of \(\eta\) by algebraic stacks is equivalent to the algebraicity of \(\mathcal X^T\). To see this, first notice that the map \(\eta:\mathcal X^T\to \mathcal X\) is always faithful. Indeed, a morphism \(\phi:(x,\{\alpha_g\})\to (y,\{\beta_g\})\) in \(\mathcal X^T\) is by definition a morphism \(\phi:x\to y\) in \(\mathcal X\) which satisfies the equivariance condition \(\phi\circ \alpha_g = \beta_g \circ (g\cdot \phi)\). This is a condition on \(\phi\), not additional data, so
  \[\Hom_{\mathcal X^T}\big((x,\{\alpha_g\}), (y,\{\beta_g\})\big) \hookrightarrow \Hom_{\mathcal X}(x,y)\]
  is injective. Next, a morphism \(f:\mathcal X\to \mathcal Y\) of algebraic stacks is faithful if and only if it is representable by algebraic spaces \cite[\href{https://stacks.math.columbia.edu/tag/04Y5}{Tag 04Y5}]{stacks-project}, and a stack admitting a morphism representable by algebraic spaces to an algebraic stack is algebraic \cite[\href{https://stacks.math.columbia.edu/tag/03YS}{Tag 03YS}]{stacks-project} hence the algebraicity of \(\mathcal X^T\) is equivalent to the representability of \(\eta\).
\end{remark}

Furthermore, the hypothesis that \(\mathcal X\) have affine diagonal cannot simply be dropped. Example \ref{example:affine-diagonal-is-necessary} is a reconstruction of an example mentioned by Romagny following~\cite[Theorem 1.2.2]{romagny_AlgebraicitySmoothnessFixed2022} found in \cite[Exp. IX, Rem. 7.4]{SGA3-2},~ illustrating that \(\mathcal X^T\) may indeed fail to be algebraic when \(\mathcal X\) does not have affine diagonal. We do not know what minimal hypotheses are necessary, but affine stabilizers are sufficient by \cite[Theorem A.20]{aranha-khan-latyntsev-etal_VirtualLocalizationRevisited2025}.
\begin{lemma}\label{rl-000F}\label{lem:two-term-truncation-of-fixed-locus-cotangent-complex-no-fixed-part}
  If \(T\) is a diagonlizable group over a field which acts on an algebraic stack of finite type over \(k\) such that \(\mathcal X^T\to \mathcal X\) is representable, then the two term truncation of \(L_{\mathcal X^T/\mathcal X}\) has no fixed part with respect to the natural \(T\)-action.
\end{lemma}

\red{This is the proof nearly verbatim from Tom's original notes. NOT COMPLETED (obviously)}
\begin{proof}
  Let \(L = \tau_{\geq -1}L_{\mathcal X^T/\mathcal X}\). First we show the vanishing \(h^0(L^f) = \Omega^f_{\mathcal X^T/\mathcal X}\), for which it suffices to check geometric points. Given \(p\in \mathcal X(k')\), the fiber of \(\mathcal X^T\) over \(p\) is the space of splittings of 
\end{proof}
\red{The following ``proof'' is something I hacked together with Claude after looking at \cite[Corollary A.31]{aranha-khan-latyntsev-etal_VirtualLocalizationRevisited2025}. It's not vetted -- it's an alternate idea which needs verification. Abandoning is also viable.}
\begin{proof}
  This may be proven directly using a Picard-stack argument, but instead we appeal to \cite[Cor. A.31]{aranha-khan-latyntsev-etal_VirtualLocalizationRevisited2025}, translating their derived stack results to classical Artin stacks.
  Let \(\varepsilon:\mathcal X^{hT}\to \mathcal X\) be the derived homotopy fixed stack \cite[Def. A.14, Rem A.15]{aranha-khan-latyntsev-etal_VirtualLocalizationRevisited2025} which is \(T\)-equivariant for the trivial action on the source. For a classical \(S\)-scheme \(U\), one has
  \begin{align*}
    \operatorname{Maps}_S(U,\mathcal X^{hT}) = \operatorname{Maps}_S^T(U, \mathcal X) = \Hom_{T-\St}(\iota(U), \mathcal X) = \mathcal X^T(U),
  \end{align*}
  so the classical truncation of \(\mathcal X^{hT}\) is \(\mathcal X^T\) and \(\eta = \varepsilon \circ \iota\) with \(\iota:\mathcal X^T\to \mathcal X^{hT}\) the classical truncation morphism. By \cite[Rem. A.19]{aranha-khan-latyntsev-etal_VirtualLocalizationRevisited2025} \(BT\) is formally proper over \(S\) and hence we may apply \cite[Cor. A.31]{aranha-khan-latyntsev-etal_VirtualLocalizationRevisited2025} to obtain \(L_{\mathcal X^{hT}/\mathcal X}\simeq (\varepsilon^*L_{\mathcal X})^{mov}[1]\), whose fixed part vanishes. Decomposition of \(D_{qc}(-\times BT)\) into weight-spaces is functorial for base change, so pullback along \(i\) commutes with the fixed/moving decomposition and thus \((i^*L_{\mathcal X^{hT}/\mathcal X})^{f} = 0\). The transitivity triangle for \(\mathcal X^T\xrightarrow{i}\mathcal X^{hT}\xrightarrow{\varepsilon}\mathcal X\) is
  \begin{align*}
    i^*L_{\mathcal X^{hT}/\mathcal X} \to L_\eta \to L_{\mathcal X^T/\mathcal X^{hT}} \xrightarrow{+1},
  \end{align*}
  which is a triangle in \(D_{qc}(\mathcal X^{T}\times_SBT)\). Since \(i\) is the inclusion of the classical truncation, \(L_{\mathcal X^T/\mathcal X^{hT}}\) is 2-connective, meaning \(h^j(\mathcal X^T/\mathcal X^{hT}) = 0\) for \(j \geq -1\). The long exact sequence in cohomology therefore gives \(h^0(i^*L_{\mathcal X^{hT}/\mathcal X}) \simeq h^0(L_\eta)\) and a surjection \(h^{-1}(L_{\mathcal X^{hT}/\mathcal X}) \twoheadrightarrow h^0(L_\eta)\). Taking fixed parts yields the claim.
\end{proof}
We end this section describing how, given an obstruction theory for \(\pi:M\to \mathcal X\), we obtain a canonical obstruction theory for \(M^T\to \mathcal X^T\).
\begin{lemma}\label{rl-000G}\label{lem:obstruction-theory-for-the-fixed-locus}
  Suppose \(E\) is a \(T\)-equivariant relative perfect obstruction theory for \(\pi:M\to \mathcal X\) where \(M\) is a Deligne-Mumford stack and \(\mathcal X\) is an algebraic stack. That is, \(E\) is an element in the \(T\)-equivariant derived category of \(M\) of perfect amplitude \([-1,0]\) together with a \(T\)-equivariant morphism \(\phi: E\to L_{M/\mathcal X}\) which is an obstruction theory in the sense of \cite[Definition 4.4]{behrend-fantechi_IntrinsicNormalCone1997}. If \(\mathcal X^T\) is an algebraic stack, then \((E|_{M^T})^f\) is a canonical relative perfect obstruction theory for \(M^T\) over \(\mathcal X^T\).
\end{lemma}
\begin{proof}
  We have a commutative square (not Cartesian)
  \begin{center}
    \begin{tikzcd}
      M^T \arrow[r, "\iota"] \arrow[d, "\pi^T"] & M \arrow[d, "\pi"] \\
      \mathcal X^T \arrow[r, "\eta"] & \mathcal X
    \end{tikzcd}.
  \end{center}
  Traversing from \(M^T\) to \(\mathcal X\) through \(M\) and through \(\mathcal X^T\) each gives us a transitivity triangle of cotangent complexes, and pulling back \(\phi\) by \(\iota^*\) gives us a morphism \(\iota^*\phi:\iota^*E\to \iota^*L_{M/\mathcal X}\). We can arrange this data as follows:
  \begin{center}
    \begin{tikzcd}
      \iota^*E \arrow[d,"\iota^*\phi "] & & & \\
    \iota^*L_{M/\mathcal X} \arrow[r,"\alpha"] & L_{M^T/\mathcal X} \arrow[r] \arrow[d, equal] & L_{M^T/M} \arrow[r, "+1"] & \dots\\
    (\pi^T)^* L_{\mathcal X^T/\mathcal X} \arrow[r] & L_{M^T/\mathcal X} \arrow[r, "\beta"] & L_{M^T/\mathcal X^T} \arrow[r, "+1"] & \dots
  \end{tikzcd}
  \end{center}
  where the middle equality follows from the commutativity of the diagram. Write \(\iota^*E\) as \(E|_{M^T}\) to match the notation in the statement and define \(\psi\) to be the composition of \(\iota^*\phi\), \(\alpha\) and \(\beta\) after taking fixed parts:
  \begin{align*}
    \psi: \left(E|_{M^T}\right)^f\xrightarrow{\iota^*\phi^f} \left(\iota^*L_{M/\mathcal X}\right)^f \xrightarrow{\alpha^f} \left(L_{M^T/\mathcal X}\right)^f \xrightarrow{\beta^f} \left(L_{M^T/\mathcal X^T}\right)^f.
  \end{align*}
  We show that \(\psi\) is a perfect obstruction theory. First consider the map \(\iota^*\phi\). Right exactness of derived pullback implies that \(h^0(\iota^*E) = \iota^*h^0(E)\) and \(h^0(\iota^*L_{M/\mathcal X}) = \iota^*h^0(L_{M/\mathcal X})\), and \(\iota^*\) preserves isomorphisms, so \(h^0(\iota^*\phi)\) is an isomorphism. For \(h^{-1}(-)\), we get a Tor correction and a diagram with exact rows:
  \begin{center}
  \begin{tikzcd}
    \iota^*h^{-1}(E) \arrow[r] \arrow[d,"\iota^*h^{-1}(\phi)"] & h^{-1}(\iota^*E) \arrow[r]\arrow[d,"h^{-1}(\iota^*\phi)"] & \Tor_1(h^0(E)) \arrow[r] \arrow[d,"\Tor_1(h^0(\phi))"] & 0 \arrow[d] \\
    \iota^*h^{-1}(L_{M/\mathcal X}) \arrow[r] & h^{-1}(\iota^*L_{M/\mathcal X}) \arrow[r] & \Tor_1(h^0(L_{M/\mathcal X})) \arrow[r] & 0 \\
  \end{tikzcd}
  \end{center}
  \(\Tor\) preserves isomorphisms so \(\Tor_1(h^0(\phi))\) is an isomorphism. The morphism \(h^{-1}(\phi)\) is surjective by assumption and \(\iota^*(-)\) is right exact, so \(\iota^*h^{-1}(\phi)\) is also surjective. Applying \cite[Lemma 05QA]{stacks-project} then gives us that \(h^{-1}\iota^*\phi\) is surjective as well.

  Now consider the map \(\alpha\). The fixed portion of the long exact sequence of the transitivity triangle involving \(\alpha\) reads
  \begin{align*}
    &h^{-1}(\iota^*L_{M/\mathcal X})^f \xrightarrow{h^{-1}(\alpha^f)}h^{-1}(L_{M^T/\mathcal X})^f \to h^{-1}(L_{M^T/M})^f \to \dots\\\dots \to ~& h^{0}(\iota^*L_{M/\mathcal X})^f \xrightarrow{h^{0}(\alpha^f)}h^{0}(L_{M^T/\mathcal X})^f \to h^{0}(L_{M^T/M})^f \to 0
  \end{align*}
  

  Since \(\iota\) is a closed immersion (Lemma \ref{lem:torus-actions-are-nice}) \(h^0(L_{M^T/M}) = \Omega_{M^T/M} = 0\), and applying Lemma \ref{lem:two-term-truncation-of-fixed-locus-cotangent-complex-no-fixed-part} to \(\mathcal M\) gives us \(h^{-1}(L_{M^T/M})^f = 0\). This implies \(h^0(\alpha^f)\) is an isormophism and \(h^{-1}(\alpha^f)\) is surjective.

  Finally consider \(\beta\). Because \(T\) acts trivially on \(\mathcal M^T\) and \(\mathcal X^T\), \(L_{\mathcal M^T/\mathcal X^T}\) is of pure weight zero and the weight-zero part of the long exact sequence of the bottom triangle reads
  \begin{align*}
    &h^{-1}(\iota^*L_{\mathcal X^T/\mathcal X})^f \to h^{-1}(L_{M^T/\mathcal X})^f \xrightarrow{h^{-1}(\beta^f)} h^{-1}(L_{M^T/\mathcal X^T})^f \to \dots\\\dots \to ~& h^{0}(\iota^*L_{\mathcal X^T/\mathcal X})^f \to h^{0}(L_{M^T/\mathcal X})^f \xrightarrow{h^{0}(\beta^f)} h^{0}(L_{M^T/\mathcal X^T})^f \to 0.
  \end{align*}
  Since \((\pi^T)^*\) is right \(t\)-exact and compatible with weights, \(h^0((\pi^T)^*L_{\mathcal X^T/\mathcal X})^f = (\pi^T)^*(\Omega^f_{\mathcal X^T/\mathcal X}) = 0\) by Lemma \ref{lem:two-term-truncation-of-fixed-locus-cotangent-complex-no-fixed-part}. Hence \(h^0(\beta^f)\) is an isomorphism and \(h^{-1}(\beta^f)\) is surjective. Thus \(h^0(\psi)\) is the composite of isomorphisms and \(h^{-1}(\psi)\) is the composite of surjections, so we conclude that \(\psi\) is a perfect obstruction theory.
\end{proof}

